% Generated draft: requires mathematical and source review.
\documentclass[11pt]{article}
\usepackage[margin=25mm]{geometry}
\usepackage{fontspec}
\setmainfont{Latin Modern Roman}
\usepackage{amsmath,amssymb,mathtools,unicode-math}
\setmathfont{Latin Modern Math}
\usepackage{enumitem,needspace,xurl,hyperref}
\setlength{\parindent}{0pt}
\setlength{\parskip}{5pt}
\setlength{\emergencystretch}{3em}
\newcommand{\N}{\mathbb{N}}
\newcommand{\Z}{\mathbb{Z}}
\newcommand{\Q}{\mathbb{Q}}
\newcommand{\R}{\mathbb{R}}
\newcommand{\C}{\mathbb{C}}
\newcommand{\F}{\mathbb{F}}
\newcommand{\A}{\mathbb{A}}
\newcommand{\PP}{\mathbb{P}}
\newcommand{\E}{\mathbb{E}}
\newcommand{\eps}{\varepsilon}
\newcommand{\dd}{\,\mathrm{d}}
\begin{document}
\section*{UnsolvedMath ID: 30001030}
\textbf{Source problem code:} catalogue entry
\subsection{Definitions and mathematical statement}
Retain the source's group notation $G_d=\Z^d\wr B$ for $d\in\{3,4\}$. Under the standard restricted-wreath convention,
\[
\Z^d\wr B=\left(\bigoplus_{b\in B}\Z^d\right)\rtimes B,
\]
where $B$ shifts the indices. This makes $\Z^d$ the lamp group and $B$ the base group. Some lamplighter discussions use the opposite order, so this convention cannot be exchanged with $B\wr\Z^d$ without confirmation. The group $B$ is not defined in the extract.

For a specified finite symmetric generating set $S$ of $G_d$, simple random walk ordinarily has uniform step law $\mu$ on $S$ and transitions $g\mapsto gs$. Its bounded harmonic functions satisfy
\[
h(g)=\sum_{s\in S}\mu(s)h(gs).
\]
To identify its Poisson boundary is to give a measured $G_d$-space $(\mathcal B,\nu)$ representing \emph{all} bounded $\mu$-harmonic functions via the Poisson transform
\[
h(g)=\int_{\mathcal B}f(g\xi)\,\dd\nu(\xi),\qquad f\in L^\infty(\mathcal B,\nu),
\]
with the usual measure-theoretic identification. Merely producing a nontrivial boundary factor or a candidate limiting configuration is not the full identification.

The source asks for this boundary in dimensions $d=3$ and $d=4$. The answer cannot be uniquely specified without $B$, the wreath-product action/order convention and the generating set or step law. The supplied assessment discusses finite-second-moment results for free metabelian groups; those hypotheses and groups do not by themselves define the written wreath-product walk. No particular lamp-configuration boundary or metabelian replacement is asserted.
\subsection{Short English statement}
Identify the full Poisson boundaries of the stated simple walks for $d=3,4$. The group $B$, wreath order and walk law are unspecified, and the supplied metabelian partial results cannot be substituted for those missing definitions.
\subsection{Sources}
\begin{itemize}
\item Supplied report: \emph{Geometric Group Theory, Hyperbolic Dynamics and Symplectic Geometry} (2009), Oberwolfach Reports, DOI 10.4171/owr/2008/42, pp. 2337--2406; proposed year 2008. \url{https://doi.org/10.4171/owr/2008/42}.
\item Supplied Lyons--Peres, \emph{JEMS}, reference: \url{https://ems.press/content/serial-article-files/32538}.
\item Supplied Erschler, \emph{Commentarii Mathematici Helvetici}, reference: \url{https://doi.org/10.4171/CMH/220}.
\item Local UnsolvedMath registry, permanent ID 30001030, catalogue entry; created 25 August 2026, updated 16 September 2026. Its corrected\_verified note says the formula was retained with notation repaired. Literature checked at \texttt{2026-08-21}: the assessment is supplied dated curation, not independently verified. It labels the record partially solved, reporting a free-metabelian boundary extension in rank at least three under finite second moment, while retaining a broader question; its relation to the exact wreath-product conventions is not specified.
\end{itemize}
\end{document}
